% FinEval v1 technical report. Every figure in this file comes from numbers.tex or tables/*.tex,
% both generated by `python -I -m fineval.report_tables`; tests/test_report.py checks that the prose
% contains no typed digits. Compile: tectonic main.tex
\documentclass[10pt,twocolumn,a4paper]{article}
\usepackage[margin=1.8cm,columnsep=0.6cm]{geometry}
\usepackage[T1]{fontenc}
\usepackage[utf8]{inputenc}
\usepackage{lmodern}
\usepackage{microtype}
\usepackage{amsmath,amssymb}
\usepackage{booktabs}
\usepackage{adjustbox}
\usepackage{xcolor}
\usepackage{enumitem}
\usepackage[numbers,sort&compress]{natbib}
\usepackage[colorlinks=true,linkcolor=blue!50!black,citecolor=blue!50!black,urlcolor=blue!50!black]{hyperref}
\usepackage{caption}
\captionsetup{font=small,labelfont=bf}
\setlist{itemsep=1pt,topsep=2pt,leftmargin=*}

\newcommand{\pending}[1]{\textcolor{red!70!black}{[pending: #1]}}
\newcommand{\fe}{\textsc{FinEval}~v1}
\makeatletter
\newcommand{\tabinput}[1]{\@@input #1 }
\makeatother
% generated by python -I -m fineval.report_tables; do not edit
\newcommand{\nAmbiguousTag}{7}
\newcommand{\nCalcTest}{1{,}000}
\newcommand{\nConflicts}{28}
\newcommand{\nFactTest}{199}
\newcommand{\nFactTrain}{2{,}014}
\newcommand{\nJoinExact}{944}
\newcommand{\nJoinId}{248}
\newcommand{\nMulti}{834}
\newcommand{\nPinned}{108}
\newcommand{\nPublishedRows}{61}
\newcommand{\nQATest}{150}
\newcommand{\nReproducedRows}{0}
\newcommand{\nReview}{6}
\newcommand{\nReviewPending}{6}
\newcommand{\nSftFincl}{24{,}798}
\newcommand{\nSftFinni}{2{,}014}
\newcommand{\nSftNoContext}{363}
\newcommand{\nSkipped}{8}
\newcommand{\nSourceDisagreeQA}{2}
\newcommand{\nSourceItems}{11{,}401}
\newcommand{\nSources}{4}
\newcommand{\nTagPrimary}{574}
\newcommand{\nTagTest}{581}
\newcommand{\nTermTest}{6{,}488}
\newcommand{\nTrainDropped}{15}
\newcommand{\nUnique}{10{,}457}
\newcommand{\retFinclAtFifty}{23.6}
\newcommand{\retFinclAtOne}{6.2}
\newcommand{\retFinclAtTwoHundred}{31.2}
\newcommand{\retFinclTypeAtFifty}{31.2}
\newcommand{\retFinclTypeAtOne}{7.6}
\newcommand{\retFinclTypeAtTwoHundred}{41.5}
\newcommand{\retFormula}{84.4}
\newcommand{\retPageTen}{34.7}
\newcommand{\retPageThree}{22.7}
\newcommand{\retTagAtFifty}{38.9}
\newcommand{\retTagAtOne}{18.9}
\newcommand{\retTagAtTwoHundred}{62.2}


\title{\textbf{\fe: One Harness for XBRL, SEC-Filing and Accounting Evaluation,\\
with Explicit Overlap Handling and Published-vs-Reproduced Scores}}
% Author order: Dev Gaglani (first / corresponding author), Stavan Desai.
% Mentor to be added once approved.
\author{Dev Gaglani\thanks{First and corresponding author.} \and Stavan Desai}
\date{\today}

\begin{document}
\twocolumn[
\maketitle
\begin{abstract}
Several public benchmarks test whether language models can read financial filings and the XBRL
taxonomy: XBRL-Agent~\cite{han2024xbrlagent}, the COLING Regulations
Challenge~\cite{wang2025regulations}, FinTagging~\cite{wang2025fintagging} and
FinanceBench~\cite{islam2023financebench}. They were built from overlapping material, so the same
question often appears in more than one of them, sometimes with a different gold answer, and their
leaderboards mix numbers measured on different item sets. We merge the four public releases into
one evaluation set with a common record format. From \nSourceItems{} source items, an
exact-or-identifier deduplication finds \nUnique{} unique items; \nMulti{} of them occur in more
than one release, \nConflicts{} carry conflicting gold answers, and near-duplicates are routed to
human review rather than merged automatically. We define five tasks (concept tagging, fact
extraction, filing QA, financial calculation and terminology), versioned scorers, a closed-book and a
tool-assisted track, and a leaderboard site that keeps scores copied from papers apart from scores we
measure. Model results with open-weight models on a laptop GPU are \pending{GPU runs}.
\end{abstract}
\vspace{1em}
]

\section{Introduction}
Evaluation of language models on financial reporting has grown quickly, but the resources overlap.
The COLING terminology questions are a subset of XBRL-Agent's; the COLING domain and numeric
questions are FinanceBench questions; and a reader comparing two leaderboard numbers usually cannot
tell whether they were measured on the same items. \fe{} is a deliberately small first version that
makes these relations explicit. It contributes:
\begin{itemize}
\item a pinned, hash-checked copy of each public source and a single record schema (\S\ref{sec:data});
\item a conservative deduplication that joins items only on a shared source identifier or identical
  normalised text, records every gold conflict, and never joins near-duplicates without a human
  decision (\S\ref{sec:dedup});
\item split rules that keep any item seen in test out of training and few-shot pools (\S\ref{sec:splits});
\item versioned scorers that follow each source's published rule where it is known (\S\ref{sec:tasks});
\item a tool-assisted track whose retrieval ceilings are reported before any model is run
  (\S\ref{sec:tools});
\item a protocol and site that show published and reproduced numbers side by side but never rank
  them together (\S\ref{sec:published}).
\end{itemize}
FinAuditing~\cite{wang2025finauditing} is out of scope for v1: its multi-document inputs do not fit
the context we can run locally, and its full data, like the full FinTagging sets, is gated on Hugging
Face.

\section{Sources}\label{sec:data}
Table~\ref{tab:sources} lists the sources. Each file is fetched from a fixed commit and stored with
its SHA-256 hash, licence and retrieval date; \nPinned{} files are pinned in total, including the
FinanceBench filings referenced by the open sample. A changed upstream file is refused, not silently
replaced. Loaders log every row they skip and why (\nSkipped{} rows, all empty XBRL-Agent terms).
FinanceBench is licensed for non-commercial use, which the datasheet and site carry forward.

\begin{table}[t]\centering\small
\caption{Sources, licences and number of pinned files.}\label{tab:sources}
\begin{adjustbox}{max width=\columnwidth}\tabinput{tables/sources.tex}\end{adjustbox}
\end{table}

Every record has a unique identifier of the form \texttt{source:subset:split:id}, the task, the input,
the gold answer, a versioned metric name, the licence, provenance and, where known, the filing it came
from. Table~\ref{tab:datasets} gives the size of each source dataset.

\begin{table}[t]\centering\small
\caption{Source datasets by task (before deduplication).}\label{tab:datasets}
\begin{adjustbox}{max width=\columnwidth}\tabinput{tables/datasets.tex}\end{adjustbox}
\end{table}

\section{Deduplication and conflicts}\label{sec:dedup}
\paragraph{Exact joins only.} Two records are joined when they share a FinanceBench identifier or
when their normalised text is identical. Normalisation applies Unicode NFKC, strips enumeration and
quotes, lowercases, canonicalises numbers and drops punctuation. Joins form clusters by union-find:
\nJoinExact{} joins come from identical text and \nJoinId{} from shared identifiers. Each cluster
keeps one representative, chosen from the original publisher first (FinanceBench, then
XBRL-Agent or FinTagging, then COLING) and from the test split first.

\paragraph{Near-duplicates go to people.} A first version that also joined fuzzy matches merged
pairs such as ``Due from'' and ``Due to'' and templated questions about different companies, so
automatic near joins were removed. Candidate pairs are found with MinHash LSH~\cite{broder1997minhash}
and a character-ratio filter, restricted to pairs from different datasets that contain the same
numbers and have no exact partner; the \nReview{} candidates are written to a review file
(\nReviewPending{} still pending) and affect nothing until a reviewer decides them.

\paragraph{Containment.} Table~\ref{tab:containment} lists every source dataset that is fully
contained in another. Such datasets add no new items; scores on them are reported under the
original dataset.

\begin{table}[t]\centering\small
\caption{Datasets fully contained in another (exact or identifier joins).}\label{tab:containment}
\begin{adjustbox}{max width=\columnwidth}\tabinput{tables/containment.tex}\end{adjustbox}
\end{table}

\paragraph{Gold conflicts.} When copies of one item disagree across sources, the original
publisher's gold is kept and the item is flagged \texttt{source\_disagreement} (\nSourceDisagreeQA{}
filing-QA items; in one, the COLING and XBRL-Agent copies of a FinanceBench working-capital question
give a positive answer where the original gives a negative one). When the same dataset contains one input with different golds, the
item is flagged \texttt{ambiguous}, keeps the union of golds and is left out of the primary score
(\nAmbiguousTag{} concept-tagging queries). Table~\ref{tab:conflicts} gives the counts.

\begin{table}[t]\centering\small
\caption{Clusters whose members carry different golds.}\label{tab:conflicts}
\begin{adjustbox}{max width=\columnwidth}\tabinput{tables/conflicts.tex}\end{adjustbox}
\end{table}

\section{Splits and leakage}\label{sec:splits}
A cluster is assigned a single split, and is test if any member is test. Training items are dropped
if they conflict, are near-duplicates of an evaluation item, or contain test text; \nTrainDropped{}
were dropped. Validation items whose filing also appears in test are flagged, and few-shot pools
exclude both the item's cluster and its filing. Table~\ref{tab:benchmark} gives the final
benchmark. The build is deterministic: rebuilding from the pinned files reproduces every output
byte for byte.

\begin{table}[t]\centering\small
\caption{\fe{} items per task and split; ``primary'' excludes ambiguous items.}\label{tab:benchmark}
\begin{adjustbox}{max width=\columnwidth}\tabinput{tables/benchmark.tex}\end{adjustbox}
\end{table}

\section{Tasks and scorers}\label{sec:tasks}
\begin{description}[leftmargin=0pt]
\item[Concept tagging] (FinTagging FinCL, COLING tag queries): map a reported value in context to its
  US-GAAP element. \texttt{concept\_exact@v1} is exact string equality, as in FinTagging's released
  scoring notebook; a case-insensitive element match is reported as a sensitivity check.
\item[Fact extraction] (FinTagging FinNI): list the numeric facts in a passage with their types.
  \texttt{finni\_pair\_f1@v1} matches (fact, type) pairs as multisets and pools precision, recall and
  F1 over items. FinTagging's own FinNI scorer is not in its public code, so this rule is ours and is
  documented as such.
\item[Filing QA] (FinanceBench open sample): \texttt{qa\_judge@v1} labels answers correct, incorrect
  or failed, as in the FinanceBench paper, with a judge validated against human labels
  (\S\ref{sec:judge}).
\item[Financial calculation] (XBRL-Agent formulas): \texttt{numeric\_match@v1} compares the final
  number with a tolerance of two decimals or a small relative error, and reads a percentage as a
  fraction when the gold is a fraction.
\item[Terminology] (XBRL-Agent and COLING terms): free-text definitions under
  \texttt{term\_judge@v1}. Its judge has not been validated yet, so terminology scores are not part of
  the v1 leaderboard.
\end{description}
Every score comes with a bootstrap confidence interval, and model comparisons use a paired bootstrap.
Missing or skipped generations score zero.

\section{Tool-assisted track}\label{sec:tools}
Besides closed-book prompting, each task has a tool track with deterministic, local tools: BM25
search~\cite{robertson2009bm25} over the US-GAAP taxonomy (optionally filtered by entity type), BM25
search over XBRL-Agent's formula list with a whitelisted arithmetic evaluator, and BM25 page search
over the FinanceBench filings. Because a tool can only help when it retrieves the answer, we report
each tool's ceiling before running any model (Table~\ref{tab:retrieval}). Taxonomy search finds the
gold element among its top two hundred candidates for \retFinclAtTwoHundred\% of FinCL queries
(\retFinclTypeAtTwoHundred\% with the type filter) and \retTagAtTwoHundred\% of COLING tag queries.
Formula search ranks the right formula first for \retFormula\% of calculation items. Page search is
weak: an evidence page is among the top three for \retPageThree\% of questions and among the top ten
for \retPageTen\%. No retrieval setting was tuned on test items.

\begin{table}[t]\centering\small
\caption{Taxonomy retrieval: gold element in top $k$ (\%).}\label{tab:retrieval}
\begin{adjustbox}{max width=\columnwidth}\tabinput{tables/retrieval.tex}\end{adjustbox}
\end{table}

\section{Published versus reproduced scores}\label{sec:published}
Published numbers are copied verbatim from the papers with the page and table they come from
(\nPublishedRows{} rows; Tables~\ref{tab:pub-fintagging}--\ref{tab:pub-coling}). Each row records
the item set it was measured on and how that set relates to ours. FinTagging's paper reports the full
gated sets, while we evaluate the public subsets, so those rows are marked as a different item set.
Published and reproduced scores are never ranked in one list: the site groups published rows by
evaluation set and shows our measured rows separately, with confidence intervals.

\begin{table}[t]\centering\small
\caption{Published FinTagging scores (\%), full sets~\cite{wang2025fintagging}.}\label{tab:pub-fintagging}
\begin{adjustbox}{max width=\columnwidth}\tabinput{tables/published_fintagging_table6.tex}\end{adjustbox}
\end{table}

\begin{table}[t]\centering\small
\caption{Published FinanceBench results (\% correct)~\cite{islam2023financebench}.}\label{tab:pub-financebench}
\begin{adjustbox}{max width=\columnwidth}\tabinput{tables/published_financebench_table2.tex}\end{adjustbox}
\end{table}

\begin{table}[t]\centering\small
\caption{Published COLING XBRL scores (\%)~\cite{wang2025regulations}.}\label{tab:pub-coling}
\begin{adjustbox}{max width=\columnwidth}\tabinput{tables/published_coling_table4.tex}\end{adjustbox}
\end{table}

\section{Results}
Closed-book and tool-assisted runs of open-weight models served locally through Ollama (Llama, Qwen
and Mistral families, quantised to fit a six-gigabyte GPU), with a second seed and a reproduction of
the FinTagging Llama baseline on the public subsets, are \pending{GPU runs}. Reproduced rows so far:
\nReproducedRows.

\section{Training}
Training data comes from FinTagging's released training set only, after the leakage rules of
\S\ref{sec:splits}: \nSftFinni{} fact-extraction examples using the evaluation instruction verbatim,
and \nSftFincl{} concept-tagging examples in the FinCL query format (\nSftNoContext{} entities whose
value could not be located in the passage were skipped). We fine-tune a small instruct model with
QLoRA~\cite{dettmers2023qlora}; held-out examples for early checks are drawn from training items only,
grouped by source item. Results are \pending{GPU runs}.

\section{Judge validation}\label{sec:judge}
Filing-QA answers are judged by a model from a different family than the model being judged. Two
annotators label a blind sample independently; we report their agreement and the judge's agreement
with their consensus. The publication threshold was fixed before any labels were seen: judge scores
are primary only if Cohen's $\kappa$ reaches a pre-set level, otherwise they are shown as
provisional. Results are \pending{GPU runs and labelling}.

\section{Website}
A static site reads only exported JSON files: leaderboard, comparison, task, methodology, dataset and
reproduction pages. A test checks that the page source contains no typed figures, and a browser check
confirms that every number rendered on the page is a formatted value from those files.

\section{Limitations}
The v1 item sets are the public subsets, not the full FinTagging or FinAuditing data. Near-duplicate
review is pending, so some paraphrased duplicates may remain as separate items. Page retrieval over
long filings is weak, which bounds the tool track for filing QA. The FinNI scorer is our own rule. All
models are small and quantised, so absolute scores are not comparable to API models in published
tables.

\bibliographystyle{plainnat}
\bibliography{refs}
\end{document}
